\chapter{Beispiel: Kapitel}\label{chap:example}
Dies ist eine Kapitel.

\section{Beispiel: Abschnitt}\label{sec:example}
Dies ist eine Abschnitt.

\section{Beispiel: Unterabschnitt}\label{subsec:example}
Dies ist eine Unterabschnitt.

\subsubsection{Beispiel: Teilabschnitt}
Dies ist ein Teilabschnitt.

\section{Beispiel: Sonstige Komponenten}
\subsection{Abbildungen}
\begin{figure}[H]
	\centering
	\fbox{\includegraphics[width=\figurewidth]{generic/img/FuBK_testcard_vectorized.eps}}
	\caption{Testbild Abbildung  \small{Quelle: \cite{Diefenbach.1961}}}
	\label{fig:Test}
\end{figure}

\subsection{Akronyme}
Beispiel von der Benutzung von Abkürzungen:\\
Bei der ersten Nutzung, werden die Abkürzungen ausgeschrieben \gls{acn:FH} \gls{acn:AC}.\\
Bei der zweiten Nutzung, nur abgekürzt: \gls{acn:FH} \gls{acn:AC}.\\
Umlautbeispiel: \gls{acn:UNB} \\

\subsection{Glossar}
\glqq In unserem Netzwerk setzen wir auf \gls{acn:AD}. Durch den Einsatz
eines \gls{acn:AD} erreichen wir bei \gls{acn:MS}-Systemen, die mit einer
\gls{glos:AntwD} von \gls{acn:CD} installiert wurden, die beste Standardisierung.

\subsection{Symbole}
Berechnungen mit \gls{symb:Pi} ergeben stets ein ungenaues Ergebnis,
denn \gls{symb:Pi} ist eine irrationale Zahl. Weiterhin gibt es noch
\gls{symb:Phi} und \gls{symb:Lambda}.

\subsection{Aufzählung}
\begin{itemize}
	\item[a)] Aufzählung A
	\item[b)] Aufzählung B
\end{itemize}
\begin{itemize}
	\item[1.] Aufzählung 1
	\item[2.] Aufzählung 2
\end{itemize}

\subsection{Formel}
\begin{figure}[H]
	\begin{equation}\label{eq:Example}
	\begin{split}
	&x = Testvariable 1\\
	&y = Testvariable 2\\
	\ \\
	f(x,y) &= 
	\begin{cases} 
	x + y \textnormal{ für } x <= 100 \\ 
	x \textnormal{ für } 100 < x  \\
	\end{cases}\\
	\end{split}
	\end{equation}
	\equations{Beispielformel \small{Quelle: [Eigene Darstellung]}}
\end{figure}

\subsection{Umlaute}
Ohne Pakete  \"A, \"O, \"U, \"a, \"o, \"u und \ss{} \\
Ohne Pakete {\"A}, {\"O}, {\"U}, {\"a}, {\"o}, {\"u} und {\ss} \\
Durch babel mit ngerman "A, "O, "U, "a, "o, "u und "s \\
Durch selinput umgewandelt Ä, Ö, Ü, ä, ö, ü und ß

\subsection{Tabelle}
\begin{table}[H]
	\centering
	\caption{Beispieltabelle 1 \small{Quelle:  [Eigene Darstellung]}}
	\label{tab:Example}
	\fbox{
		\spacebetween{1.5}
		\begin{tabular}{c|c|c}
			\textbf{Überschrift 1} & \textbf{Überschrift 2} & \textbf{Überschrift 3}\\
			\hline
			\hline
			Viel Inhalt & Viel Inhalt & Viel Inhalt\\
			\hline
			Viel Inhalt & Viel Inhalt & Viel Inhalt\\
			\hline
			Viel Inhalt & Viel Inhalt & Viel Inhalt\\
		\end{tabular}
	}
\end{table}


\begin{table}[H]
	\centering
	\caption{Beispieltabelle 2 \small{Quelle:  [Eigene Darstellung]}}
	\resizebox{\tablewidth}{!}{
		\spacebetween{1.3}
		\begin{tabular}{@{}rrrrcrrrcrrr@{}}
			\toprule
			& \multicolumn{3}{c}{$w = 8$} & \phantom{abc}& \multicolumn{3}{c}{$w = 16$} &
			\phantom{abc} & \multicolumn{3}{c}{$w = 32$}\\
			\cmidrule{2-4} \cmidrule{6-8} \cmidrule{10-12}
			& $t=0$ & $t=1$ & $t=2$ && $t=0$ & $t=1$ & $t=2$ && $t=0$ & $t=1$ & $t=2$\\ \midrule
			$dir=1$\\
			$c$ & 0.0790 & 0.1692 & 0.2945 && 0.3670 & 0.7187 & 3.1815 && -1.0032 & -1.7104 & -21.7969\\
			$c$ & -0.8651& 50.0476& 5.9384&& -9.0714& 297.0923& 46.2143&& 4.3590& 34.5809& 76.9167\\
			$c$ & 124.2756& -50.9612& -14.2721&& 128.2265& -630.5455& -381.0930&& -121.0518& -137.1210& -220.2500\\
			$dir=0$\\
			$c$ & 0.0357& 1.2473& 0.2119&& 0.3593& -0.2755& 2.1764&& -1.2998& -3.8202& -1.2784\\
			$c$ & -17.9048& -37.1111& 8.8591&& -30.7381& -9.5952& -3.0000&& -11.1631& -5.7108& -15.6728\\
			$c$ & 105.5518& 232.1160& -94.7351&& 100.2497& 141.2778& -259.7326&& 52.5745& 10.1098& -140.2130\\
			\bottomrule
		\end{tabular}
	}
\end{table}

\begin{table}[H]
	\centering
	\caption{Beispieltabelle 3 \small{Quelle: [Eigene Darstellung]}}
	\spacebetween{1}
	\begin{tabular}{@{}cccc@{}}
		\toprule
		procesy & Min & Max & Delta \\
		\midrule
		Junior & 1,700 & 3,200 &  1,500  \\
		\multicolumn{4}{@{}c@{}}{\dashrule} \\
		Senior & ? & ? & ? \\
		\multicolumn{4}{@{}c@{}}{\dashrule} \\
		Specialist & ? & ? & ? \\
		\multicolumn{4}{@{}c@{}}{\dashrule} \\
		Team Leader & ? & ? & ? \\
		\bottomrule
	\end{tabular}
\end{table}
	

\subsection{Text}
\tiny{tiny}\\ % smallest size
\scriptsize{scriptsize}\\ % size of supscripts
\footnotesize{footnotesize}\\ % size of footnotes
\small{small}\\
\normalsize{normalsize}\\ % size of regular text
\large{large}\\
\Large{Large}\\
\LARGE{LARGE}\\
\huge{huge}\\
\Huge{Huge}\\ % biggest size

\subsection{Code}

\begin{figure}[h]
\begin{lstlisting}[language=python]
def hello():
    print("Hello World!")
\end{lstlisting}
\caption{Beispielprogrammcode}
\end{figure}

\subsection{Zitierweisen}
\normalsize

\subsubsection{DIN 1505 Standard}
\glqq Das Wort zitieren geht zurück auf lateinisch citare, das >>herbeirufen<<, auch >>anrufen, erwähnen, nennen<< bedeutet.\grqq [Dud17]

\subsubsection{APA-Stil}
\glqq Das Wort zitieren geht zurück auf lateinisch citare, das >>herbeirufen<<, auch >>anrufen, erwähnen, nennen<< bedeutet.\grqq [Dudenredaktion.2017]

\subsubsection{Deutscher-Stil}
\glqq Das Wort zitieren geht zurück auf lateinisch citare, das >>herbeirufen<<, auch >>anrufen, erwähnen, nennen<< bedeutet.\grqq\footnote{Dudenredaktion (Hg.) (2017): Zitate und Aussprüche. Herkunft und aktueller Gebrauch. 4. Auflage. Berlin: Bibliographisches Institut; Duden (Duden - Deutsche Sprache in 12 Bänden, 12).}

\subsection{Weiter Beispiele}
http://www.texample.net/\\
http://www.texample.net/tikz/examples/feature/trees/